\documentclass[11pt]{article}
\usepackage[margin=1in]{geometry}
\usepackage{amsmath,amssymb,amsthm,mathrsfs}
\usepackage[T1]{fontenc}
\usepackage[hidelinks]{hyperref}
\usepackage{microtype}
\setlength{\parindent}{0pt}
\setlength{\parskip}{0.62em}
\setlength{\emergencystretch}{3em}

\newtheorem{theorem}{Theorem}
\newtheorem{lemma}[theorem]{Lemma}
\newcommand{\cZ}{\mathcal Z}
\newcommand{\R}{\mathbb R}
\newcommand{\Z}{\mathbb Z}

\title{Five Zak Zeros for a Generalized First Hermite Function\\
\large A counterexample to Conjecture 4.1 of arXiv:2502.09510}
\author{fa\_banach\_001, lane 4}
\date{August 17, 2026}

\begin{document}
\maketitle

\begin{abstract}
Faulhuber conjectured that every dilated and chirped first Hermite function
has exactly the three Zak-transform zeros forced by odd parity.  We give an
exact hexagonal counterexample.  For
\[
 a=\sqrt{\frac{2}{\sqrt3}},\qquad s=\frac12,
\]
the generalized first Hermite function has two additional Zak zeros at
\((1/6,5/6)\) and \((5/6,1/6)\).  The proof identifies the squared norm of a
Gaussian Zak transform with an absolutely convergent theta series and uses
its order-three affine symmetry.  No numerical approximation enters the
proof.
\end{abstract}

\section{Source conjecture and result}

For a suitable function \(f\), use the convention
\[
 \cZ f(x,\omega)=\sum_{k\in\Z}f(k-x)e^{2\pi i k\omega}.
 \tag{1}
\]
The first Hermite function is a nonzero constant multiple of
\(t e^{-\pi t^2}\).  Section 4 of M.~Faulhuber,
\emph{Gabor systems with Hermite functions of order \(n\) and oversampling
greater than \(n+1\) which are not frames}, arXiv:2502.09510, makes the
following conjecture.  If
\[
 \widetilde h_1(t;a,s)
 =\frac1{\sqrt a}h_1\!\left(\frac ta\right)e^{\pi i s t^2},
 \qquad a>0,\quad s\in\R,
 \tag{2}
\]
then \(\cZ\widetilde h_1\) has exactly three zeros in \([0,1)^2\), namely
\[
 (0,0),\qquad (1/2,0),\qquad (0,1/2).
 \tag{3}
\]
These three zeros are indeed always present because (2) is odd.

\begin{theorem}\label{thm:main}
Conjecture 4.1 of arXiv:2502.09510 is false.  Set
\[
 a_0=\sqrt{\frac{2}{\sqrt3}},\qquad s_0=\frac12.
 \tag{4}
\]
In addition to the three zeros in \emph{(3)}, one has
\[
 \cZ\widetilde h_1\!\left(\frac16,\frac56;a_0,s_0\right)=0,
 \qquad
 \cZ\widetilde h_1\!\left(\frac56,\frac16;a_0,s_0\right)=0.
 \tag{5}
\]
Thus this generalized first Hermite function has at least five distinct Zak
zeros in the fundamental square.
\end{theorem}

\section{Gaussian reduction}

Let
\[
 q=a^{-2}-is,\qquad \tau=iq=s+\frac{i}{a^2},\qquad
 \phi_q(t)=e^{-\pi q t^2},
 \tag{6}
\]
so that \(\Re q>0\) and \(\Im\tau>0\).  Put
\[
 \Theta(x,\omega)=\cZ\phi_q(x,\omega),\qquad
 F(x,\omega)=|\Theta(x,\omega)|^2.
 \tag{7}
\]
All series below, together with their derivatives, converge normally.

\begin{lemma}\label{lem:critical}
At every point where \(\Theta\ne0\),
\[
 \nabla F=0\quad\Longleftrightarrow\quad \Theta_x=0.
 \tag{8}
\]
Moreover, for a nonzero constant \(C_{a,s}\),
\[
 \cZ\widetilde h_1(x,\omega;a,s)=C_{a,s}\Theta_x(x,\omega).
 \tag{9}
\]
\end{lemma}

\begin{proof}
Since \(h_1(t)=C_1t e^{-\pi t^2}\) with \(C_1\ne0\), (2) and (6) give
\[
 \widetilde h_1(t;a,s)=C_1a^{-3/2}t\,e^{-\pi q t^2}.
\]
Termwise differentiation of (7) yields
\[
 \Theta_x
 =2\pi q\sum_{k\in\Z}(k-x)e^{-\pi q(k-x)^2}e^{2\pi i k\omega},
\]
which proves (9).

The same series gives the first-order identity
\[
 \Theta_\omega=\frac{i}{q}\Theta_x+2\pi i x\Theta.
 \tag{10}
\]
At a point where \(\Theta\ne0\), write \(r=\Theta_x/\Theta\).  The equation
\(F_x=0\) says \(\Re r=0\), hence \(r=it\) for some \(t\in\R\).  By (10),
\[
 \frac{F_\omega}{2F}
 =\Re\!\left(\frac{i}{q}r+2\pi i x\right)
 =-t\,\Re\!\left(\frac1q\right).
\]
Because \(\Re q>0\), one has \(\Re(1/q)>0\); therefore \(F_\omega=0\) if
and only if \(t=0\), or equivalently \(\Theta_x=0\).  The reverse
implication follows from the same formulas.
\end{proof}

\section{The hexagonal affine symmetry}

Write
\[
 \tau=c+ib,\qquad b>0,
\]
so \(q=b-ic\).  The next Fourier expansion is the source of the exact
symmetry.

\begin{lemma}\label{lem:fourier}
The periodic function \(F\) has the absolutely convergent Fourier series
\[
 F(x,\omega)
 =\frac1{\sqrt{2b}}\sum_{n,m\in\Z}
 (-1)^{mn}
 \exp\!\left(-\frac{\pi}{2b}|n+m\tau|^2\right)
 e^{2\pi i(nx+m\omega)}.
 \tag{11}
\]
\end{lemma}

\begin{proof}
Let \(\widehat F(n,m)\) denote the Fourier coefficient obtained by
multiplying by \(e^{-2\pi i(nx+m\omega)}\) and integrating over the unit
square.  Expanding \(|\Theta|^2\), integration in \(\omega\) imposes a
difference of indices equal to \(m\).  Joining the remaining unit intervals
then gives
\[
 \widehat F(n,m)
 =\int_\R \phi_q(t+m)\overline{\phi_q(t)}e^{2\pi i nt}\,dt.
\]
Completing the square in this Gaussian integral yields
\[
 \widehat F(n,m)
 =\frac1{\sqrt{2b}}
 \exp\!\left(-\frac{\pi}{2b}|n+m\tau|^2-\pi i mn\right),
\]
which is (11).  Gaussian decay also proves absolute convergence.
\end{proof}

For the parameters (4),
\[
 \tau=\frac12+i\frac{\sqrt3}{2},\qquad b=\frac{\sqrt3}{2},
\qquad \tau^2=\tau-1.
 \tag{12}
\]
Define the integer matrix and its inverse transpose by
\[
 A=\begin{pmatrix}0&-1\\1&1\end{pmatrix},
 \qquad
 B=A^{-T}=\begin{pmatrix}1&-1\\1&0\end{pmatrix}.
 \tag{13}
\]
The map \(A\) sends \((n,m)\) to \((-m,n+m)\), and (12) gives
\[
 \tau(n+m\tau)=-m+(n+m)\tau.
 \tag{14}
\]
It therefore preserves the Gaussian factor in (11).  Its sign changes by
\[
 (-1)^{(-m)(n+m)}=(-1)^m(-1)^{mn}.
 \tag{15}
\]
Reindexing (11) with (14)--(15) proves
\[
 F\bigl(S(x,\omega)\bigr)=F(x,\omega),\qquad
 S(x,\omega)=\left(x-\omega+\frac12,x\right)\pmod{\Z^2}.
 \tag{16}
\]

\begin{lemma}\label{lem:fixed}
The points
\[
 p_1=\left(\frac16,\frac56\right),\quad
 p_0=\left(\frac12,\frac12\right),\quad
 p_2=\left(\frac56,\frac16\right)
 \tag{17}
\]
are critical points of \(F\).
\end{lemma}

\begin{proof}
From (13) and (16),
\[
 S^2(x,\omega)
 =\left(-\omega,x-\omega+\frac12\right)\pmod{\Z^2}.
 \tag{18}
\]
Solving \(S^2(x,\omega)=(x,\omega)\) on the torus gives exactly the three
points (17).  The derivative of \(S^2\) is
\[
 B^2=\begin{pmatrix}0&-1\\1&-1\end{pmatrix},
\]
whose eigenvalues are the primitive cube roots of unity; in particular it
has no eigenvalue \(1\).  At a fixed point \(p\), differentiating
\(F\circ S^2=F\) gives
\[
 (B^2)^T\nabla F(p)=\nabla F(p).
\]
Thus \(\nabla F(p)=0\).
\end{proof}

\section{Proof of the counterexample}

It remains only to exclude the possibility that the critical points
\(p_1,p_2\) are zeros of \(\Theta\), because Lemma \ref{lem:critical} was
stated away from those zeros.  With the standard Jacobi theta function
\[
 \vartheta_3(z;\tau)
 =\sum_{k\in\Z}e^{\pi i\tau k^2+2\pi i k z},
\]
expansion of the square in (7) gives
\[
 \Theta(x,\omega)
 =e^{\pi i\tau x^2}\,
   \vartheta_3(\omega-\tau x;\tau).
 \tag{19}
\]
The classical theta function \(\vartheta_3(\,\cdot\,;\tau)\) has one simple
zero modulo \(\Z+\tau\Z\), at \((1+\tau)/2\).  Consequently \(\Theta\) has
exactly one zero modulo \(\Z^2\), namely
\[
 (x,\omega)=\left(\frac12,\frac12\right)=p_0.
 \tag{20}
\]
For completeness, the theta zero count follows immediately by integrating
\(\vartheta_3'/\vartheta_3\) around a period parallelogram and using the
standard quasi-periodicity; the displayed half-period is a zero by pairing
terms.

It follows from (20) that \(\Theta(p_1)\) and \(\Theta(p_2)\) are nonzero.
Lemma \ref{lem:fixed} and Lemma \ref{lem:critical} now give
\[
 \Theta_x(p_1)=\Theta_x(p_2)=0.
\]
Equation (9) proves the two identities (5).  Together with the three
parity-forced zeros (3), this completes the proof of Theorem
\ref{thm:main}.

\section{Scope, verification, and novelty}

The theorem is a complete negative answer to the paper's explicit Conjecture
4.1.  It does not determine the full zero set, nor does it answer Questions
Q1--Q3 for every Hermite order.  Applying the source paper's periodic-zero
criterion to the five displayed zeros also gives a density-five periodic
non-frame Gabor system for this generalized first Hermite window.

The accompanying high-precision script evaluates all five zeros and checks
the affine norm symmetry at sample points.  It is only a regression test; the
proof above is exact.

A bounded novelty search covered the run indexes, the exact conjecture,
the two rational zero locations, generalized first Hermite functions,
hexagonal lattices, and related theta/Zak terminology.  It also inspected
the closely related arXiv:2502.08358, which uses a theta identity for a
second-Hermite example.  No prior statement of (5) or answer to Conjecture
4.1 was found.  Long--Wei--Zou's 2025 classification of critical points of
heat kernels on tori reports the closely analogous six-critical-point
phenomenon on the hexagonal torus, providing conceptual corroboration but
not, in the bounded search, this Zak/Hermite specialization.  Novelty is
therefore provisional pending expert bibliographic review.

\begin{thebibliography}{9}
\bibitem{Faulhuber}
M. Faulhuber,
\emph{Gabor systems with Hermite functions of order \(n\) and oversampling
greater than \(n+1\) which are not frames},
Sampling Theory, Signal Processing, and Data Analysis 23 (2025);
arXiv:2502.09510.

\bibitem{FaulhuberZak}
M. Faulhuber,
\emph{Operator-isomorphism pairs and Zak transform methods for the study of
Gabor systems},
arXiv:2502.08358, 2025.

\bibitem{LongWeiZou}
C.-G. Long, J. Wei, and W. Zou,
\emph{The classification of critical points of heat kernel on tori and
application to elliptic equation and Mueller--Ho conjecture},
Ann. Sc. Norm. Super. Pisa Cl. Sci., 2025.
\end{thebibliography}

\end{document}
